\Needspace{22\baselineskip}
\Article{4}
I now turn to \emph{specific} arrangements of the quarter systems, and set the modulus, to be denoted by $m$, $=ea+ebi$, so that the positive integer $e$ denotes the greatest real divisor, or the greatest divisor common to the two components of $m$; thus $a,b$ are relatively prime. The quarter system with the simplest arrangement will be the one whose members $x+iy$ are such that $ax+by$ is positive and less than $\frac12e(aa+bb)$, while $ay-bx$ is nonnegative and likewise less than $\frac12e(aa+bb)$; the latter condition is deliberately expressed so as to include the cases in which $ay-bx=0$. It is easy to see that there will be $\frac12(e-1)$ such cases in all, namely

\[
\begin{array}{rlrl}
x&=a,&y&=b\\
x&=2a,&y&=2b\\
x&=3a,&y&=3b\\
\multicolumn{4}{c}{\text{and so on, up to}}\\
x&=\frac12(e-1)a,&y&=\frac12(e-1)b
\end{array}
\]
and hence none at all if the components of $m$ have no common divisor.
\SourceLeaf{382}
Let us call this quarter system $k$, and let $k',k'',k'''$ be those obtained by multiplying the numbers belonging to $k$ by $i,-1,-i$; or one may also put
\[
k'=m+ik,\quad k''=(1+i)m-k,\quad k'''=im-ik
\]

This gives the following rule for deciding whether any given integer $x+iy$ not divisible by $m$ is congruent to a member of $k,k',k''$ or $k'''$. Namely, $2(ax+by)$ can be put in the form $Pe(aa+bb)+Q$, and $2(ay-bx)$ in the form $Re(aa+bb)+S$, with $P,Q,R,S$ real integers satisfying the following conditions:

\begin{center}if $x+iy$ is congruent to a number from\end{center}
\begingroup\small\setlength{\arraycolsep}{3pt}
\[
\begin{array}{l||l|l|l|l}
\text{so that}&k&k'&k''&k'''\\
P&\text{even}&\text{odd}&\text{odd}&\text{even}\\
R&\text{even}&\text{even}&\text{odd}&\text{odd}\\
Q&\text{positive}&\text{positive}&\text{positive}&\text{nonnegative}\\
e(aa+bb)-Q&\text{positive}&\text{nonnegative}&\text{positive}&\text{positive}\\
S&\text{nonnegative}&\text{positive}&\text{positive}&\text{positive}\\
e(aa+bb)-S&\text{positive}&\text{positive}&\text{nonnegative}&\text{positive}
\end{array}
\]
\endgroup

The matter is easier to survey if the cases in which neither of the numbers $ax+by$, $ay-bx$ is divisible by $e(aa+bb)$ are distinguished from the rest.

I. In the first case, simply take for $P$ the (algebraically) smaller of the two integers between which $\frac{2(ax+by)}{e(aa+bb)}$ lies strictly, and likewise for $R$ the smaller of the two between which $\frac{2(ay-bx)}{e(aa+bb)}$ lies.

II. If $ax+by$ is divisible by $e(aa+bb)$, then $ay-bx$ will be divisible by $aa+bb$, but not by $e(aa+bb)$ (since otherwise $x+iy$ would be divisible by $ea+ebi$). If now $R$, that is, the integer immediately below $\frac{2(ay-bx)}{e(aa+bb)}$, is even, then $x+iy$ will be congruent modulo $ea+ebi$ to a number from $k'$, but to one from $k'''$ if $R$ is odd.

III. If $ay+bx$ is divisible by $e(aa+bb)$, but $ax+by$ is not, then $x+iy$ will be congruent to a number from $k$ or from $k''$, according as the integer algebraically immediately below $\frac{2(ax+by)}{e(aa+bb)}$ is even or odd.

\SourceLeaf{383}

\Article{5}
The following method for determining the character of a given integer $M=A+Bi$ relative to the odd modulus $m=ea+ebi$ that does not divide it is readily derived from the above.

For brevity we use the following notation. If $p$ denotes any nonintegral real number, let $[p]$ denote the integer that makes both $p-[p]$ and $1+[p]-p$ positive. This definition therefore excludes applying the notation to integers. If the definition required neither $p-[p]$ nor $1+[p]-p$ to be negative, the symbol $[p]$ would be ambiguous when $p$ is an integer. One could also, as was done in an earlier paper, require $1+[p]-p$ to be positive and $p-[p]$ merely to be nonnegative. For our purpose it is somewhat more convenient to keep to the first definition.

Accordingly, the quarter system $k$ consists of all integers $f$ for which, on putting $\frac{2f}{m}=\xi+\eta i$, $\xi$ lies strictly between $0$ and $1$, while $\eta$ lies between $0$ and $1$, with $0$ included; or
\[
\begin{array}{lll}
&[\xi]=0,&[\eta]=0\\
\text{or}&[\xi]=0,&\eta=0
\end{array}
\]

If now, for every $f$, we put $\frac{2fM}{m}=\xi+i\eta$ and take
\begingroup\small\setlength{\arraycolsep}{2pt}
\[
\begin{array}{rl}
\Psi f=0&\text{if both }[\xi]\text{ even and either }[\eta]\text{ even or }\eta\text{ integer}\\
1&[\eta]\text{ even and either }[\xi]\text{ odd or }\xi\text{ integer}\\
2&[\xi]\text{ odd and either }[\eta]\text{ odd or }\eta\text{ integer}\\
3&[\eta]\text{ odd and either }[\xi]\text{ even or }\xi\text{ integer}
\end{array}
\]
\endgroup

or, in tabular form,
\[
\begin{array}{l|c|c|c}
&[\eta]\text{ even}&[\eta]\text{ odd}&\eta\text{ integer}\\\hline
[\xi]\text{ even}&0&3&0\\
\relax[\xi]\text{ odd}&1&2&2\\
\xi\text{ integer}&1&3&-\ \pdot
\end{array}
\]
which may be denoted by $\nabla(\xi+i\eta)$, then the required character of the number $M$ relative to the modulus $m$ equals the least residue of $\Sigma\Psi f$ modulo $4$.

\SourceLeaf{384}

\Needspace{18\baselineskip}
\Article{6}
The rule given in the preceding article is general; for the case in which $M$ is odd, however, we shall give it another form. At the same time we shall assume that the real parts of $m$ and $M$ are odd, and hence their imaginary parts even.

To each number $f$ belonging to the quarter system $k$ we assign another number $g$, derived from $f$ as follows. On putting $\frac{2f}{m}=\xi+i\eta$, four cases are distinguished:
\[
\begin{array}{ll}
\mathrm{I.}&\text{If }[2\xi]=0\text{ and either }[2\eta]=0\text{ or }\eta=0\\
\mathrm{II.}&\text{If }[2\xi]=1\text{ and either }[2\eta]=0\text{ or }\eta=0\\
\mathrm{III.}&\text{If }[2\xi]=0\text{ and }[2\eta]=1\\
\mathrm{IV.}&\text{If }[2\xi]=1\text{ and }[2\eta]=1
\end{array}
\]
In the first case put $g=2f$, in the second $g=im-2if$, in the third $g=m+2if$, and in the fourth $g=(1+i)m-2f$. It is easy to see that the collection of all $g$ forms a complete quarter system $l$; its defining property is that, on putting $\frac gm=\xi+i\eta$, the following conditions hold together:

\[
\begin{array}{llll}
\text{either}&[\xi]=0,&[\eta]=0&\\
\text{or}&\eta=0,&[\xi]=0&\text{and }g\text{ divisible by }1+i\\
\text{or}&\xi=0,&[\eta]=0&\text{and }g\text{ not divisible by }1+i
\end{array}
\]
It follows that $l$ differs from $k$ only in that the numbers in $k$ for which $\eta=0$ and which are not divisible by $1+i$, namely
\[
\begin{aligned}
&a+bi,\quad3(a+bi),\quad5(a+bi)\ldots\frac{e-3}{2}(a+bi)\\
&\text{or up to }\frac{e-1}{2}(a+bi)
\end{aligned}
\]
according as $e$ has the form $4n+1$ or $4n+3$, are absent from $l$, while their products by $i$ occur in the latter collection instead. At the same time one sees that, for $e=1$, that is, when $m$ is divisible by no real integer, $k$ and $l$ are entirely the same.

It now remains to derive $\Psi f$ directly from the $g$ corresponding to $f$. The result is that, for the above $4$ cases,
\[
\begin{aligned}
\mathrm{I.}\quad&\Psi f=\nabla\frac{gM}{m}\\
\mathrm{II.}\quad&\Psi f=\left\{\begin{array}{ll}
-\nabla\frac{gM}{m}&\begin{gathered}\text{if neither the real nor imaginary part}\\\text{ of }\frac{gM}{m}\text{ is an integer}\end{gathered}\\
1-\nabla\frac{gM}{m}&\text{if either is an integer}
\end{array}\right.
\end{aligned}
\]
III. as II.\quad IV. as I.
\BookRule{.3\textwidth}
